% §1 Exercises
%
% Cross-references and their resolution status:
%
% Already resolved (within chapter 1):
%   \Cref{sec:three-faces}                  -- §1.1
%   \Cref{sec:quantifying-noise}            -- §1.2
%   \Cref{sec:boolean-universe-first-look}  -- §1.3
%   \Cref{sec:sampling-rates}               -- §1.4
%   \Cref{sec:why-atom}                     -- §1.5
%
% Forward references (resolved later in the book or in spirit only):
%   Chapter 16 (negative-space) is named in prose only (no \Cref);
%   chapter 16 is not yet drafted, and the references in exercises 2,
%   10, and 11 are intentionally prose-level forward gestures.
%   Chapter 7 (Wald CI alternatives) is named in prose only.
%
% Difficulty-marker macros (not in alex.sty; defined locally if not yet defined):
\providecommand{\exstar}{$\star$}
\providecommand{\exstarstar}{$\star\star$}
\providecommand{\exstarstarstar}{$\star\star\star$}

\section*{Exercises}
\addcontentsline{toc}{section}{Exercises}

\begin{enumerate}

% -----------------------------------------------------------------------
% Difficulty tiers (markers, not physical order): the exercises stay in
% their natural teaching order; the star markers carry the difficulty.
%   Warm-up      (★   × 3): exercises 1, 2, 3
%   Medium       (★★  × 5): exercises 4, 5, 6, 8, 9
%   Challenging  (★★★ × 3): exercises 7, 10, 11
% -----------------------------------------------------------------------
% WARM-UP  (★ × 3): exercises 1, 2, 3
% -----------------------------------------------------------------------

\item \exstar
\label{ex:ch1-1}
\textbf{Confusion matrix to rate conversion.}
A binary classifier is evaluated on a labelled test set and produces the
following counts: $\mathrm{TP} = 90$, $\mathrm{FP} = 10$, $\mathrm{FN} = 10$,
$\mathrm{TN} = 90$.
\begin{enumerate}
  \item[(a)] Compute the empirical false-positive rate $\hat{\fprate}$
    and false-negative rate $\hat{\fnrate}$ using the row-conditional
    proportions from \Cref{sec:quantifying-noise}. Be explicit about
    which counts go in the numerator and which in the denominator.
  \item[(b)] Compute the precision $\mathrm{TP}/(\mathrm{TP}+\mathrm{FP})$
    and the recall $\mathrm{TP}/(\mathrm{TP}+\mathrm{FN})$.
  \item[(c)] In one sentence, explain why $\hat{\fprate}$ and precision are
    \emph{different} numbers even though both involve the false-positive
    count. (\textit{Hint:} which marginal does each one condition on?)
\end{enumerate}

\item \exstar
\label{ex:ch1-2}
\textbf{Three real-world noisy bits and one non-example.}
The three faces of \Cref{sec:three-faces} (BSC, biased coin, classifier
verdict) are stand-ins for many real systems. Name \emph{three} systems
of your own that you think of as noisy bits in the sense of
\Cref{sec:quantifying-noise} (a Boolean-output system with stable
per-input rates).
\begin{enumerate}
  \item[(a)] For each of the three, identify what plays the role of the
    Boolean input and Boolean output, then write a rough estimate of
    $\fprate$ and $\fnrate$.
  \item[(b)] Construct one explicit \emph{non-example} whose per-input
    error rates are not stable: they change over time, or depend on
    which queries came before. The error behaviour is real, but the
    rates are not the stable numbers the noisy-bit model of
    \Cref{sec:quantifying-noise} assumes.
  \item[(c)] In one sentence, identify the stability property of the
    noisy-bit model that your non-example violates. We return to
    non-examples like this in chapter 16, which makes the failure
    modes of the model its explicit subject.
\end{enumerate}

\item \exstar
\label{ex:ch1-3}
\textbf{Joint table to rates.}
A noisy bit is described by the joint table below, where the truth has
prior $\Pr\{X = \top\} = 0.3$ and $\Pr\{X = \bot\} = 0.7$:
\[
  \begin{array}{c|cc}
    & Y = \top & Y = \bot \\\hline
    X = \top & 0.27 & 0.03 \\
    X = \bot & 0.07 & 0.63
  \end{array}
\]
The four cell entries are joint probabilities $\Pr\{X, Y\}$, and they sum
to 1.
\begin{enumerate}
  \item[(a)] Verify that the row sums match the priors.
  \item[(b)] Compute the row-conditional false-positive rate $\fprate =
    \Pr\{Y = \top \Given X = \bot\}$ and the false-negative rate
    $\fnrate = \Pr\{Y = \bot \Given X = \top\}$ by dividing the
    appropriate cell by its row sum (\Cref{sec:quantifying-noise}).
  \item[(c)] Confirm that $\fprate$ and $\fnrate$ are properties of the
    noisy bit itself, not of the prior $\Pr\{X = \top\}$: recompute
    them for a different prior (say $\Pr\{X = \top\} = 0.5$) with the
    same conditional behaviour, and check that they are unchanged. In
    one sentence, explain why this independence from the prior is what
    lets the noisy-bit model get by with two rates.
\end{enumerate}

% -----------------------------------------------------------------------
% MEDIUM  (★★ × 5): exercises 4, 5, 6, 8, 9
% (exercise 7 sits in the natural order here but is marked ★★★)
% -----------------------------------------------------------------------

\item \exstarstar
\label{ex:ch1-4}
\textbf{Empirical estimation under varying $n$ (notebook-tied).}
Open the chapter 1 companion notebook
\texttt{python/notebooks/ch01-noisy-bit.ipynb}. Use the binary symmetric
channel simulator with true crossover $\fprate = 0.10$ and run the
estimator at $n \in \{10, 100, 1000, 10000\}$.
\begin{enumerate}
  \item[(a)] Tabulate $\hat{\fprate}$ for each value of $n$ from a single
    run, and the absolute deviation $|\hat{\fprate} - \fprate|$ from
    the true rate.
  \item[(b)] Plot $\hat{\fprate}$ as a function of $n$ on a log-$n$
    axis, with the true rate drawn as a horizontal reference line.
    Describe in one sentence how the spread of the estimator around
    the true rate shrinks with $n$.
  \item[(c)] Compare the rate at which the empirical deviation shrinks
    to the predicted standard error
    $\sqrt{\fprate(1-\fprate)/n}$ from \Cref{sec:sampling-rates}.
    Is the observed convergence consistent with a $1/\sqrt{n}$ shrinking
    band?
\end{enumerate}

\item \exstarstar
\label{ex:ch1-5}
\textbf{The trivial biased coin.}
A trivial biased coin is a noisy bit that ignores the input entirely:
it returns $\top$ with probability $p$ and $\bot$ with probability
$1-p$, regardless of whether the truth is $\top$ or $\bot$.
\begin{enumerate}
  \item[(a)] Using the definitions in \Cref{sec:quantifying-noise},
    write $\fprate$ and $\fnrate$ as functions of $p$.
  \item[(b)] At what value of $p$ is the trivial biased coin
    \emph{useless} for the classification task, in the sense that
    $\fprate$ and $\fnrate$ are both equal? Compute that pair of rates.
  \item[(c)] At what values of $p$ does the trivial biased coin achieve
    its \emph{best} behaviour, in the sense that one of the two rates
    is exactly zero? Which rate, at which value of $p$?
  \item[(d)] Comment briefly on the asymmetry: a constant function is
    a perfect classifier for one truth class and a maximally bad one
    for the other. The biased-coin parameterization makes this trade-off
    visible.
\end{enumerate}

\item \exstarstar
\label{ex:ch1-6}
\textbf{NOT swaps the two rates.}
Let $b$ be a noisy bit with rates $(\fprate, \fnrate)$. Define $b' = \neg
b$ to be the noisy bit obtained by swapping the output: whenever $b$
reports $\top$, $b'$ reports $\bot$, and vice versa. (No additional
randomness is introduced; $b'$ is a deterministic relabelling of the
output of $b$.)
\begin{enumerate}
  \item[(a)] Write out the four conditional probabilities
    $\Pr\{Y = \top \Given X = \top\}$, $\Pr\{Y = \bot \Given X = \top\}$,
    $\Pr\{Y = \top \Given X = \bot\}$, $\Pr\{Y = \bot \Given X = \bot\}$
    for $b$, following the row-conditional convention of
    \Cref{sec:quantifying-noise}.
  \item[(b)] Apply the output swap to $b$ and write out the same four
    conditional probabilities for $b'$.
  \item[(c)] Read off the rates of $b'$. Show that $\fprate' = \fnrate$
    and $\fnrate' = \fprate$.
  \item[(d)] Conclude that NOT acts on the parameter pair as the swap
    $(\fprate, \fnrate) \mapsto (\fnrate, \fprate)$. In one sentence,
    interpret this geometrically as a reflection in the parameter plane
    $[0,1]^2$.
\end{enumerate}

\item \exstarstarstar
\label{ex:ch1-7}
\textbf{Threshold sweep and the error-rate trade-off curve.}
A score-based classifier produces a score $s(x) \in [0, 1]$ on each
input $x$ and reports $\top$ if and only if $s(x) > \theta$ for a
threshold $\theta$. For each value of $\theta$, the system is a noisy
bit with rates $\fprate(\theta)$ and $\fnrate(\theta)$ that depend on
$\theta$.
\begin{enumerate}
  \item[(a)] As $\theta$ increases from $0$ to $1$, describe what
    happens to $\fprate(\theta)$ and $\fnrate(\theta)$: which one
    increases, which one decreases, and why. (You do not need a
    formula; reason from the threshold rule of \Cref{sec:quantifying-noise}.)
  \item[(b)] At $\theta = 0$ the classifier reports $\top$ for almost
    every input; at $\theta = 1$ it reports $\bot$ for almost every
    input. Give the $(\fprate, \fnrate)$ pair at each of these two
    extremes.
  \item[(c)] The curve traced by $(\fprate(\theta), \fnrate(\theta))$
    as $\theta$ sweeps from $0$ to $1$ is the trade-off curve between
    the two error rates. The standard \emph{ROC curve} plots the
    true-positive rate against the false-positive rate, which is the
    same information re-axed. Sketch (in words) the shape of the
    $(\fprate, \fnrate)$ trade-off curve, and explain why a classifier
    cannot in general drive both rates to zero at once.
\end{enumerate}

\item \exstarstar
\label{ex:ch1-8}
\textbf{Sample-size formula via the CLT bound.}
\Cref{sec:sampling-rates} gives the variance of the empirical
false-positive estimator as $\Var(\hat{\fprate}) = \fprate(1 - \fprate)/n_0$.
Suppose we want the standard error to satisfy $\sqrt{\Var(\hat{\fprate})}
\le \sigma$ for some target $\sigma > 0$.
\begin{enumerate}
  \item[(a)] Solve the inequality $\sqrt{\fprate(1 - \fprate)/n_0} \le
    \sigma$ for $n_0$. Write the smallest integer $n_0$ that satisfies
    the bound as a function of $\fprate$ and $\sigma$.
  \item[(b)] Evaluate $n_0$ for the practical target $\sigma = 0.01$ at
    $\fprate \in \{0.01, 0.05, 0.10, 0.50\}$. Tabulate the four
    sample-size requirements.
  \item[(c)] Identify the value of $\fprate$ that requires the largest
    sample size. In one sentence, explain why this maximum sits at
    $\fprate = 1/2$ rather than at the boundary (cf.\
    \Cref{sec:sampling-rates}, which observes that the variance numerator
    peaks at $1/2$).
\end{enumerate}

% -----------------------------------------------------------------------
% CHALLENGING / EXTENSION  (★★★ × 3): exercises 7, 10, 11
% (exercise 9 sits in the natural order here but is marked ★★)
% -----------------------------------------------------------------------

\item \exstarstar
\label{ex:ch1-9}
\textbf{Informal LLN argument for $\hat{\fprate} \to \fprate$.}
Argue informally that under iid sampling, the empirical estimator
$\hat{\fprate}$ from \Cref{sec:sampling-rates} converges to the true
rate $\fprate$ in probability as $n_0 \to \infty$.
\begin{enumerate}
  \item[(a)] State precisely the iid assumption you are using: each of
    the $n_0$ negative-truth queries produces a Bernoulli$(\fprate)$
    output independent of all the others. Write the empirical count
    $\textnormal{FP}$ as a sum of $n_0$ such Bernoulli variables.
  \item[(b)] Recall (or look up) Chebyshev's inequality: for any random
    variable $Z$ with finite variance and any $\delta > 0$,
    $\Pr\{|Z - \mathbb{E}[Z]| \ge \delta\} \le \Var(Z)/\delta^2$.
    Apply this to $Z = \hat{\fprate}$, using the unbiasedness statement
    $\mathbb{E}[\hat{\fprate}] = \fprate$ and the variance formula
    $\Var(\hat{\fprate}) = \fprate(1 - \fprate)/n_0$ from
    \Cref{sec:sampling-rates}.
  \item[(c)] Show that the right-hand side of the Chebyshev bound goes
    to zero as $n_0 \to \infty$ for any fixed $\delta > 0$. Conclude
    that $\hat{\fprate} \to \fprate$ in probability.
  \item[(d)] In one sentence, identify which assumption from part (a)
    is most likely to fail in a real data system, and what kind of
    correlation across queries would prevent the Chebyshev argument
    from going through.
\end{enumerate}

\item \exstarstarstar
\label{ex:ch1-10}
\textbf{Stable rates, correlated errors.}
The noisy-bit model rests on two assumptions: the per-input rates are
stable across queries, and the noise events on different queries are
independent. Construct an explicit Boolean-input, Boolean-output system
with \emph{stable} per-input rates whose errors are nonetheless
\emph{correlated} across queries, so that a naive iid estimator
misjudges its behaviour.
\begin{enumerate}
  \item[(a)] Construct the system. Specify the input domain (Boolean),
    the output domain (Boolean), and the mechanism that determines the
    output given the input. The mechanism must produce a well-defined
    Boolean output on every query and must have well-defined, stable
    per-input rates $\fprate$ and $\fnrate$, but the error events on
    successive queries must be statistically dependent in a stated way.
  \item[(b)] Identify precisely how the independence assumption fails:
    describe the dependence between the error event on one query and
    the error event on the next, and confirm in one sentence that the
    per-input rates themselves are still stable (so it is independence,
    not stability, that breaks).
  \item[(c)] Argue that a naive observer who computes $\hat{\fprate}$
    and $\hat{\fnrate}$ from a sample of queries, and then treats those
    estimates as if the queries were iid, would misjudge the system. Be
    specific: the point estimates may be fine, but what quantity (an
    error bar, a coverage claim, a sample-size requirement) would they
    get wrong, and in which direction?
  \item[(d)] Chapter 16 is the negative-space chapter of the book: its
    explicit subject is what fails to fit the model. Your construction
    is a candidate worked example for that chapter; in one sentence,
    state the conceptual lesson your example would carry there.
\end{enumerate}

\item \exstarstarstar
\label{ex:ch1-11}
\textbf{Stress test of the notebook coverage estimator.}
The chapter 1 notebook \texttt{python/notebooks/ch01-noisy-bit.ipynb}
runs a Wald-interval coverage check that empirically counts the
fraction of $95\%$ intervals containing the true $\fprate$. \Cref{sec:sampling-rates} flags the under-coverage caveat:
the Wald form is asymptotically correct but is known to under-cover at
small $n_0$ when $\fprate$ sits near $0$ or $1$.
\begin{enumerate}
  \item[(a)] Identify parameter regimes (combinations of $n_0$ and
    $\fprate$) under which you predict the empirical coverage will fall
    visibly below $0.95$. Justify your prediction by reference to the
    under-coverage caveat in \Cref{sec:sampling-rates}.
  \item[(b)] Run the notebook's coverage estimator in two regimes of
    your choosing: one where you predict good coverage (close to $0.95$)
    and one where you predict bad coverage. Report the two empirical
    coverage numbers.
  \item[(c)] Explain in one paragraph \emph{why} the Wald form
    under-covers in the bad regime. The argument should reference the
    Normal approximation underlying the Wald form and identify what
    fails when $\fprate$ sits near a boundary or $n_0$ is small. (No
    rigorous derivation is needed; the explanation should be sound.)
  \item[(d)] State which of the alternative intervals named in
    \Cref{sec:sampling-rates} (Wilson, Clopper-Pearson, Jeffreys) you
    would expect to perform better in your bad-coverage regime, and
    briefly explain the intuition. Full derivations are deferred to
    chapter 7.
\end{enumerate}

\end{enumerate}
